\FloatBarrier
\section{FT and Moller-cone interface clearances}
\label{sec:ft-joints}

This October 16 revision corrects the remaining FT/Moller-cone interfaces after the individual mesh studies.
Five meshes receive local material-removing operations: the inner extension, lead fill, lead case, main cone,
and tungsten add-on. FTInnerShield, FTFlange, and the cone tip remain byte-identical. Placements, passive
rotations, materials, names, and run keys are unchanged. All SQLite banks remain in \texttt{default}, indexed
by run number.

Several original engineering bores are coarse inscribed polygons. Their facets intrude into neighboring
surfaces even when nominal vertex radii agree. The inner extension also has a small shoulder at the main-cone
joint. The corrections use closed, polygonal cutters whose faces enclose the desired circular clearance
profiles, together with a conservative convex envelope of the actual lead fill. Cutters are intermediate
construction objects and are not additional detector volumes.

The inner-extension bore clears the placed FT shield, which is offset -369.73 mm in the common cone frame. Its
end is relieved around the unchanged flange bore. The lead-fill bore clears the inner extension. A scaled
convex envelope of the pre-clearance lead fill supplies at least 0.02002 mm plane clearance before final
casing simplification; a separate local cutter clears the unchanged flange. This envelope reduces casing
material loss from 4.4702\% in a rejected circular-envelope trial to 1.2093\%. The main cone receives a short
stepped relief around the inner-extension shoulder and a profiled opening around the cone tip. Its upper face
moves from z=712.002991 to 711.950012 mm. The add-on bore clears the main cone by a nominal 0.10 mm radial
profile margin. The health report retains every cutter profile and input hash.

All five corrected meshes retain one closed manifold component and their previous genera, with no detected
self-intersections or zero-area facets. The existing beam passages and radial ports remain open. A 0.0001 mm
weld of the saved float32 meshes removes numerical fragments. The casing is simplified from its boolean
reference to 16000 triangles with endpoint collapses and a maximum sampled deviation of 0.007648 mm; its
volume changes by +0.00002023\% during this simplification. It contains more triangles than the earlier
10000-triangle casing but remains below the engineering source. The eight-part total changes from 59540 to
60100 triangles (+0.9405\%). The material changes in Table~\ref{tab:ft-joints-meshes} are intentional
clearance removals, separate from the earlier decimation volume guards.

An initial cut end face passed overlap checks but produced nine GeomNav1002 boundary-navigation warnings
during tracking. Retriangulating that annular face as an angular strip retains its exact surface, boundary
coordinates, triangle count, and material volume. The previously failing electron seed passes after this
change. The initial benchmark is retained as an excluded trial and is not used in the final timing table.
Canonical STL output changes facet order and cyclic starting vertices without changing geometry or winding;
all eight outputs and the casing reference reproduce byte for byte.

Independent placed-surface checks find zero triangle contacts and zero contained vertices in both directions
at all eight corrected interfaces. Table~\ref{tab:ft-joints-gaps} gives minimum vertex-to-surface samples, not
proven global distance bounds. Three seeds (12345--12347) with 20000-point Geant4 overlap checks produce no
exceptions. The previous eight-part fixture produced 21 GeomVol1002 blocks. Matched meridian sections show the
local changes in Figures~\ref{fig:ft-joints-1} and~\ref{fig:ft-joints-2}; the FT shield uses its actual
placement offset.

The isolated tracking comparison includes the corrected eight-part assembly in run 5874 with \texttt{default}.
Matched 5 GeV electrons, protons, positive pions, and neutrons originate at the target, uniform in solid angle
over 0.8--4.8 degrees and full azimuth. FTFP\_BERT uses one thread, default production cuts, no field or
optical processes, and no GUI or event output. Five alternating-order repetitions use 200 events per particle
and seeds 12345--12349. Executable and primary hashes are retained. Table~\ref{tab:ft-joints-tracking} reports
median Geant4 beamOn wall times excluding initialization, before/corrected ratios, and paired ratio ranges.
These timings include changed clearances and material, so they are not a measurement of decimation alone.

Median tracking ratios are 0.8568, 0.8541, 0.8465, 0.8920 for electrons, protons, positive pions, and
neutrons, respectively. The corrected median tracking times are 12.1--18.1\% longer than the pre-clearance
baseline. Faster-pair counts are 0/5, 0/5, 0/5, 0/5 in the same order. All 52 final benchmark processes
complete. The 26 corrected processes report no Geant4 exceptions. The baseline overlap diagnostic reports 21
GeomVol1002 blocks. The raw per-seed results retain process CPU, wall time, initialization, and memory
separately from transport.

Loading process wall time changes from 1.279579 to 1.272204 s; CPU changes from 1.190 to 1.230 s. Startup peak
RSS is 180.453 versus 180.188 MiB. The earlier individual-piece benchmarks remain historical measurements of
their stated geometries.

All 9 affected complete run configurations pass construction without Geant4 exceptions. All 17 beamline
configurations pass volume-count construction checks, and the comparator passes 15 native geometry, 15
material, and 14 CAD comparisons. Direct database checks preserve all GEMC2 tables and all GEMC3 physical
rows; every other STL retains its previous hash in all three runtime trees. The historical add-on optimization
remains reproducible using the archived main-cone dependency. These geometry and isolated timing checks do not
establish detector-response equivalence or full CLAS12 performance.

\begin{table}[htbp]
\centering
\caption{Triangle counts and volume changes relative to the previous optimized meshes.}
\label{tab:ft-joints-meshes}
\setlength{\tabcolsep}{5pt}
\begin{tabular}{lrrr}
\toprule
Mesh & Before & Corrected & Volume change (\%) \\
\midrule
Inner extension & 12500 & 11200 & -0.689274 \\
Lead case & 10000 & 16000 & -1.209323 \\
Lead fill & 8000 & 3958 & -0.192478 \\
Tungsten add-on & 12000 & 8356 & -0.374681 \\
Main cone & 5000 & 8546 & -0.118539 \\
\bottomrule
\end{tabular}
\end{table}

\begin{table}[htbp]
\centering
\caption{Placed interface checks; distances sample vertices against the other surface.}
\label{tab:ft-joints-gaps}
\setlength{\tabcolsep}{5pt}
\begin{tabular}{lrr}
\toprule
Interface & Sampled gap (mm) & Contacts \\
\midrule
Shield / inner extension & 0.062045 & 0 \\
Flange / inner extension & 0.053650 & 0 \\
Flange / case & 0.020020 & 0 \\
Case / lead fill & 0.019957 & 0 \\
Lead fill / inner extension & 0.053503 & 0 \\
Main cone / inner extension & 0.054603 & 0 \\
Cone tip / main cone & 0.052979 & 0 \\
Add-on / main cone & 0.094198 & 0 \\
\bottomrule
\end{tabular}
\end{table}

\begin{table}[htbp]
\centering
\caption{Eight-part tracking; five repetitions of 200 events per particle.}
\label{tab:ft-joints-tracking}
\setlength{\tabcolsep}{5pt}
\begin{tabular}{lrrrr}
\toprule
Particle & Before (s) & After (s) & Ratio & Paired range \\
\midrule
Electron & 19.36220 & 22.59820 & 0.8568 & 0.8543--0.8773 \\
Proton & 5.29483 & 6.19904 & 0.8541 & 0.8507--0.8899 \\
Positive pion & 6.94979 & 8.21029 & 0.8465 & 0.8001--0.8923 \\
Neutron & 5.55687 & 6.22962 & 0.8920 & 0.8433--0.9094 \\
\bottomrule
\end{tabular}
\end{table}

\FloatBarrier
\begin{figure}[p]
\centering
\includegraphics[width=\linewidth,height=.80\textheight,keepaspectratio]{figures/ft_joint_sections_1.pdf}
\caption{Matched meridian sections before and after the local clearance corrections. Blue and orange identify
the two
named parts. Coordinates use the common passive cone frame; the shield is translated by its actual offset.
Solid outlines are sections of the stored STLs, with matched axes within each pair.}
\label{fig:ft-joints-1}
\end{figure}
\begin{figure}[p]
\centering
\includegraphics[width=\linewidth,height=.80\textheight,keepaspectratio]{figures/ft_joint_sections_2.pdf}
\caption{Matched meridian sections before and after the local clearance corrections. Blue and orange identify
the two
named parts. Coordinates use the common passive cone frame; the shield is translated by its actual offset.
Solid outlines are sections of the stored STLs, with matched axes within each pair.}
\label{fig:ft-joints-2}
\end{figure}
